\subsection{The limit of a sequence}

A sequence is the simplest setting in which the idea of a limit can be made precise. The input does not move continuously: it takes the values
\[
1,2,3,\ldots
\]
and we ask what happens to the terms far along the sequence.

\begin{definitionbox}[Limit of a sequence]
We write
\[
\lim_{n\to\infty}a_n=L
\]
when, for every \(\varepsilon>0\), there exists an index \(N\in\mathbb N\) such that
\[
n\ge N \quad\Longrightarrow\quad |a_n-L|<\varepsilon.
\]
\end{definitionbox}

The definition says more than ``the terms get close to \(L\).'' It says that after some index they remain inside every prescribed interval
\[
(L-\varepsilon,L+\varepsilon).
\]
The first terms may behave irregularly; convergence concerns the tail of the sequence.

\begin{examplebox}[The sequence \(1/n\)]
Let
\[
a_n=\frac1n.
\]
To prove that \(a_n\to0\), let \(\varepsilon>0\). We need
\[
\left|\frac1n-0\right|=\frac1n<\varepsilon.
\]
It is enough to choose an integer \(N>1/\varepsilon\). Then every \(n\ge N\) satisfies \(1/n<\varepsilon\). Hence
\[
\lim_{n\to\infty}\frac1n=0.
\]
\end{examplebox}

\begin{examplebox}[Dominant powers]
For
\[
a_n=\frac{2n+1}{n+3},
\]
divide numerator and denominator by \(n\):
\[
a_n=\frac{2+1/n}{1+3/n}.
\]
Since \(1/n\to0\), the numerator approaches \(2\) and the denominator approaches \(1\). Therefore
\[
\lim_{n\to\infty}\frac{2n+1}{n+3}=2.
\]
\end{examplebox}

\begin{examplebox}[Oscillation prevents convergence]
The sequence
\[
b_n=(-1)^n
\]
alternates between \(-1\) and \(1\). Its terms do not eventually remain close to one number, so the sequence has no limit.
\end{examplebox}

\begin{theorembox}[Basic consequences of convergence]
Every convergent sequence has exactly one limit and is bounded.
\end{theorembox}

These statements are structural. Uniqueness makes the notation \(\lim a_n\) meaningful, while boundedness says that a convergent sequence cannot escape arbitrarily far away.